\begin{scanpage}{0108}{91}
\[\left.\begin{aligned}\mathbf E&=-\frac1{ik}(\operatorname{grad}\operatorname{div}\mathbf A+k^2\mathbf A),\\\mathbf H&=\operatorname{rot}\mathbf A\end{aligned}\right\}\tag{16.01}\]
we find the following expressions for the fringing field in the wave zone:
\[\left.\begin{aligned}E_x=H_y&=ikA_x,\\E_y=-H_x&=ikA_y\end{aligned}\right\}\quad\text{with }\vartheta=0\tag{16.02}\]
and
\[\left.\begin{aligned}E_x=-H_y&=ikA_x,\\E_y=H_x&=ikA_y\end{aligned}\right\}\quad\text{with }\vartheta=\pi.\tag{16.03}\]
The vector potential is determined by the equation
\[\begin{aligned}\mathbf A={}&\frac1c\frac{e^{ikr}}r\int_0^{2\pi}\left[\int_0^{l_1}\mathbf j_1(\zeta)e^{\pm ik\zeta\cos\omega}(a-\zeta\sin\omega)d\zeta\right.\\&\left.\quad+\int_0^{l_2}\mathbf j_2(\zeta)e^{\mp ik\zeta\cos\Omega}(a-\zeta\sin\Omega)d\zeta\right]d\psi.\end{aligned}\tag{16.04}\]
Here $r$ is the distance from the discontinuity to the observation point, $\mathbf j_1(\zeta)$ is the surface current density flowing on the irradiated side of the body, and $\mathbf j_2(\zeta)$ is the current density on the shadowed side. The upper sign in the exponents refers to the case $\vartheta=\pi$, and the lower sign refers to the case $\vartheta=0$. Since the nonuniform part of the current is concentrated mainly in the vicinity of the discontinuity, the vector potential corresponding to it may be represented in the form
\[\begin{aligned}\mathbf A={}&\frac ac\frac{e^{ikr}}r\int_0^{2\pi}\left(\int_0^\infty\mathbf j_1^1(\zeta)e^{\pm ik\zeta\cos\omega}d\zeta\right.\\&\left.\quad+\int_0^\infty\mathbf j_2^1(\zeta)e^{\mp ik\zeta\cos\Omega}d\zeta\right)d\psi.\end{aligned}\tag{16.05}\]
Obviously the nonuniform part of the current near the discontinuity of a conical surface may be considered to be approximately the same as on a corresponding wedge (Figure 29). In the local cylindrical coordinate system $r_1$, $\phi_1$, $z_1$, the field from the nonuniform part
\end{scanpage}
